\begin{ex}
Đường gấp khúc ABC gồm bao nhiêu đoạn thẳng?
\shortans{2}
\loigiai{
Gồm 2 đoạn thẳng là AB và BC.
}
\end{ex}

\begin{ex}
Đường gấp khúc ABCD gồm bao nhiêu đoạn thẳng?
\shortans{3}
\loigiai{
Gồm 3 đoạn thẳng: AB, BC, CD.
}
\end{ex}

\begin{ex}
Đường gấp khúc MNPD gồm 3 đoạn thẳng có độ dài lần lượt là 3 cm, 4 cm, 5 cm. Tính độ dài đường gấp khúc MNPD (đơn vị cm).
\shortans{12}
\loigiai{
Độ dài đường gấp khúc: $3 + 4 + 5 = 12$ cm.
}
\end{ex}

\begin{ex}
Đường gấp khúc ABC có AB dài 12 cm, BC dài 15 cm. Tính độ dài đường gấp khúc ABC (đơn vị cm).
\shortans{27}
\loigiai{
Độ dài: $12 + 15 = 27$ cm.
}
\end{ex}

\begin{ex}
Đường gấp khúc gồm 4 đoạn thẳng bằng nhau, mỗi đoạn dài 5 cm. Tính độ dài đường gấp khúc đó (đơn vị cm).
\shortans{20}
\loigiai{
Độ dài: $5 \times 4 = 20$ cm.
}
\end{ex}

\begin{ex}
Đường gấp khúc dài 30 cm gồm 3 đoạn thẳng, biết hai đoạn đầu dài lần lượt là 8 cm và 12 cm. Đoạn thứ ba dài bao nhiêu xăng-ti-mét?
\shortans{10}
\loigiai{
Đoạn thứ ba dài: $30 - 8 - 12 = 10$ cm.
}
\end{ex}

\begin{ex}
Một hình tứ giác có bao nhiêu cạnh?
\shortans{4}
\loigiai{
Hình tứ giác có đúng 4 cạnh.
}
\end{ex}

\begin{ex}
Một hình tứ giác có bao nhiêu đỉnh?
\shortans{4}
\loigiai{
Hình tứ giác có đúng 4 đỉnh.
}
\end{ex}

\begin{ex}
Một hình tứ giác có độ dài các cạnh lần lượt là 3 cm, 4 cm, 5 cm, 6 cm. Tính chu vi hình tứ giác đó (đơn vị cm).
\shortans{18}
\loigiai{
Chu vi: $3 + 4 + 5 + 6 = 18$ cm.
}
\end{ex}

\begin{ex}
Một hình vuông có cạnh dài 5 cm. Chu vi hình vuông đó là bao nhiêu xăng-ti-mét?
\shortans{20}
\loigiai{
Chu vi: $5 \times 4 = 20$ cm.
}
\end{ex}

\begin{ex}
Một hình chữ nhật có chiều dài 6 cm, chiều rộng 4 cm. Chu vi hình chữ nhật đó là bao nhiêu xăng-ti-mét?
\shortans{20}
\loigiai{
Chu vi: $(6 + 4) \times 2 = 20$ cm.
}
\end{ex}

\begin{ex}
Đường gấp khúc khép kín gồm 4 đoạn thẳng tạo thành hình gì có 4 cạnh? Hình đó có bao nhiêu đỉnh?
\shortans{4}
\loigiai{
Đó là hình tứ giác có 4 đỉnh.
}
\end{ex}

\begin{ex}
Đường gấp khúc khép kín gồm 3 đoạn thẳng tạo thành hình gì? Hình tam giác có bao nhiêu cạnh?
\shortans{3}
\loigiai{
Hình tam giác có 3 cạnh.
}
\end{ex}

\begin{ex}
Một sợi dây đồng uốn thành hình tứ giác có độ dài các cạnh là 10 cm, 12 cm, 8 cm, 15 cm. Chiều dài sợi dây đồng là bao nhiêu xăng-ti-mét?
\shortans{45}
\loigiai{
Chiều dài sợi dây: $10 + 12 + 8 + 15 = 45$ cm.
}
\end{ex}

\begin{ex}
Một đường gấp khúc gồm 2 đoạn thẳng dài 24 cm, đoạn thứ nhất dài 10 cm. Đoạn thứ hai dài bao nhiêu xăng-ti-mét?
\shortans{14}
\loigiai{
Đoạn thứ hai dài: $24 - 10 = 14$ cm.
}
\end{ex}

\begin{ex}
Một hình tứ giác có 4 cạnh bằng nhau, mỗi cạnh dài 6 cm. Chu vi hình tứ giác đó là bao nhiêu xăng-ti-mét?
\shortans{24}
\loigiai{
Chu vi: $6 \times 4 = 24$ cm.
}
\end{ex}

\begin{ex}
Một con ốc sên bò theo đường gấp khúc gồm 3 đoạn: đoạn 1 dài 15 cm, đoạn 2 dài 20 cm, đoạn 3 dài 25 cm. Con ốc sên đã bò được quãng đường dài bao nhiêu xăng-ti-mét?
\shortans{60}
\loigiai{
Quãng đường: $15 + 20 + 25 = 60$ cm.
}
\end{ex}

\begin{ex}
Chu vi một hình tam giác là 24 cm. Nếu 3 cạnh bằng nhau thì mỗi cạnh dài bao nhiêu xăng-ti-mét?
\shortans{8}
\loigiai{
Mỗi cạnh dài: $24 : 3 = 8$ cm.
}
\end{ex}

\begin{ex}
Chu vi một hình tứ giác là 36 cm, tổng độ dài ba cạnh là 27 cm. Cạnh thứ tư dài bao nhiêu xăng-ti-mét?
\shortans{9}
\loigiai{
Cạnh thứ tư dài: $36 - 27 = 9$ cm.
}
\end{ex}

\begin{ex}
Có bao nhiêu hình tứ giác trong một hình vuông được chia làm 4 ô vuông nhỏ bởi 2 đường thẳng chữ thập?
\shortans{9}
\loigiai{
Có 4 hình vuông nhỏ, 4 hình chữ nhật ghép đôi, và 1 hình vuông lớn. Tổng cộng có 9 hình tứ giác.
}
\end{ex}

\begin{ex}
Một đường gấp khúc có 5 đoạn thẳng bằng nhau, mỗi đoạn dài 4 cm. Độ dài đường gấp khúc là bao nhiêu xăng-ti-mét?
\shortans{20}
\loigiai{
Độ dài: $4 \times 5 = 20$ cm.
}
\end{ex}

\begin{ex}
Một khung ảnh hình chữ nhật có chu vi là 30 cm, nửa chu vi khung ảnh là bao nhiêu xăng-ti-mét?
\shortans{15}
\loigiai{
Nửa chu vi: $30 : 2 = 15$ cm.
}
\end{ex}

\begin{ex}
Đường gấp khúc ABCDE gồm bao nhiêu đoạn thẳng?
\shortans{4}
\loigiai{
Gồm 4 đoạn thẳng: AB, BC, CD, DE.
}
\end{ex}

\begin{ex}
Đường gấp khúc đó có bao nhiêu điểm đỉnh?
\shortans{5}
\loigiai{
Có 5 điểm đỉnh: A, B, C, D, E.
}
\end{ex}

\begin{ex}
Tính độ dài đường gấp khúc gồm 3 đoạn thẳng dài lần lượt 14 cm, 16 cm và 20 cm.
\shortans{50}
\loigiai{
Độ dài: $14 + 16 + 20 = 50$ cm.
}
\end{ex}
